\documentclass[11pt]{article}
\usepackage[a4paper,margin=26mm]{geometry}
\usepackage{amsmath,amssymb,amsthm}
\usepackage[T1]{fontenc}
\usepackage{lmodern}
\usepackage[hidelinks]{hyperref}
\newtheorem{proposition}{Proposition}
\newcommand{\tmix}{t_{\mathrm{mix}}}
\title{The Johnson Family $J(N,1)$ Has No\newline Total-Variation Cutoff}
\author{Conjecture 00000003961}
\date{}
\begin{document}
\maketitle
\begin{abstract}
The assertion of total-variation cutoff for all $k\leq N/2$ is false. For the valid family $k=1$, the Johnson graph is $K_N$. Its standard half-lazy random walk has $\tmix(1/4)=2$ and $\tmix(3/4)=1$ for every $N\geq5$. Their ratio is constantly $2$, contrary to the cutoff criterion that it tend to $1$. We prove the exact transition and total-variation formulas, not merely an asymptotic estimate.
\end{abstract}

\section*{Definitions and the claimed range}
The vertices of $J(N,1)$ are the singleton subsets of $\{1,\ldots,N\}$. Distinct singletons have symmetric difference of size $2$, so $i\mapsto\{i\}$ identifies this actual Johnson graph with the complete graph $K_N$. We use the standard lazy walk: stay at the current vertex with probability $1/2$, and otherwise choose uniformly among its $N-1$ neighbors. Thus
\[
 P(i,j)=\begin{cases}1/2,&i=j,\\[2pt]1/[2(N-1)],&i\ne j.\end{cases}
 \qquad \pi(j)=\frac1N.
\]
The matrix is positive and doubly stochastic, so the uniform distribution $\pi$ is stationary. In particular, there is no disconnectedness or periodicity obstruction in this example.

Write
\[
 d_N(t)=\max_i\frac12\sum_j|P^t(i,j)-\pi(j)|,
 \qquad \tmix^{(N)}(\varepsilon)=\min\{t\geq0:d_N(t)\leq\varepsilon\}.
\]
The standard total-variation cutoff criterion is that, for every $0<\varepsilon<1/2$,
\begin{equation}\label{eq:cutoff}
 \frac{\tmix^{(N)}(\varepsilon)}{\tmix^{(N)}(1-\varepsilon)}\longrightarrow1
 \quad (N\longrightarrow\infty).
\end{equation}
See~\cite{HermonPeres} for this formulation. The source quantifies over all $k\leq N/2$ and does not require $k\to\infty$. Hence the family $k=1$, $N\geq5$ is within its stated range.

\begin{proposition}
For the half-lazy walk on $J(N,1)$, $N\geq5$, set $a_N=(N-2)/[2(N-1)]$. For every nonnegative integer $t$,
\begin{equation}\label{eq:tv}
 d_N(t)=\frac{N-1}{N}a_N^t.
\end{equation}
In particular, $\tmix^{(N)}(1/4)=2$ and $\tmix^{(N)}(3/4)=1$.
\end{proposition}
\begin{proof}
Let $U$ be the $N\times N$ matrix all of whose entries equal $1/N$. Direct multiplication gives $U^2=U$. The transition probabilities and induction then give
\[
 P=a_N I+(1-a_N)U,\qquad P^t=a_N^t I+(1-a_N^t)U.
\]
For a fixed starting state $i$, subtracting $\pi$ gives the diagonal difference $a_N^t(1-1/N)$ and each of the other $N-1$ differences $-a_N^t/N$. Since $a_N\geq0$, their absolute values sum to $2(N-1)a_N^t/N$. This proves~\eqref{eq:tv}; it is independent of $i$, so the maximum over initial states is exactly the same value. Also $a_N<1$, so these distances tend to zero for each fixed $N$ and every positive tolerance has a finite mixing time.

At the first three times, the formula becomes
\[
 d_N(0)=\frac{N-1}{N}>\frac34,\qquad
 d_N(1)=\frac{N-2}{2N},\qquad
 d_N(2)=\frac{(N-2)^2}{4N(N-1)}.
\]
For $N\geq5$, $d_N(1)>1/4$ because $N>4$, while $d_N(1)<1/2<3/4$. Furthermore,
\[
 (N-2)^2\leq N(N-1)\quad\Longleftrightarrow\quad4\leq3N,
\]
so $d_N(2)\leq1/4$. These comparisons show both the upper bounds and the minimality of the asserted mixing times.
\end{proof}

Taking $\varepsilon=1/4$ in~\eqref{eq:cutoff} now yields the constant ratio $2/1=2$, whose limit is not $1$. Thus this family has no total-variation cutoff. This already refutes the universal cutoff clause and hence the conjecture as written; no analysis of its proposed cutoff-time multiplier is needed. This elementary complete-graph calculation is not claimed as a new result about Markov chains. The conclusion does not address a different statement restricted to growing $k$.

\section*{Formal verification and reproduction}
The accompanying Lean project uses Lean 4.19.0 and Mathlib at commit
\begin{center}\small\texttt{c44e0c8ee63ca166450922a373c7409c5d26b00b}.\end{center}
It indexes $N\geq5$ by $N=n+5$. The actual singleton-subset graph, bijective labeling, adjacency, and neighbor count are verified. The transition matrix is matched to that graph. Mathlib's \texttt{doublyStochastic} structure verifies its probability entries and every matrix power, and the uniform distribution is proved stationary.

The matrix-power identity, exact finite total-variation sum, convergence, and existence of mixing times are formalized. Mixing time is the infimum of the actual nonempty set of admissible integer times; its least values $2$ and $1$ are proved. Universal quantification over starting states implements the maximum bound. The final theorem \texttt{no\_total\_variation\_cutoff} contradicts the real-limit criterion~\eqref{eq:cutoff}. No spectral-gap assertion is assumed.

The exact \href{https://github.com/The-Last-Math-Competition/The-Last-Math-Competition/blob/main/conjectures/00000003961.md}{repository statement} is in \texttt{SOURCE.md}; build commands and detailed correspondence are in \texttt{README.md}. Dependencies use public, pinned Git revisions. The verification directory records actual fresh builds, axiom audits, hashes, and PDF checks. The auxiliary \texttt{verify.py} constructs the subset graphs and multiplies rational matrices at five finite sizes. It supplements the general Lean proof.

\begin{thebibliography}{1}
\bibitem{HermonPeres}
J. Hermon and Y. Peres, \emph{On sensitivity of mixing times and cutoff},
Electronic Journal of Probability \textbf{23} (2018), paper 25.
\href{https://arxiv.org/abs/1610.04357}{arXiv:1610.04357}.
The cutoff ratio criterion appears in the abstract; the counterexample calculation here is self-contained.
\end{thebibliography}
\end{document}
